\honest{The Beauty of Settled Science}{Eliezer Yudkowsky}{2008}

Facts need not be unexplained to be beautiful, truths are worth learning even if others know them, and beliefs are worth holding even if many share them. So if you care only about controversial science, ``you will end up with a head stuffed full of garbage.''

The media report only the frontier. You never see the headlines ``General Relativity still governing planetary orbits'' or ``Phlogiston theory remains false.'' By the time science is solid, it is no longer news. Newsworthy science ``is often based on the thinnest of evidence and wrong half the time.'' I give no source. \nb{Later evidence fits the claim. Dumas-Mallet and colleagues (2017) found that ``Only 48.7\% of the 156 studies reported by newspapers were confirmed by the corresponding meta-analyses.''}

Controversies are problems hard enough that people who have spent years mastering the field can still fool themselves, which is what makes the heated arguments the media like. For outsiders they are not even fun, except for the fun of picking a side, which you can get at any football game.

A good textbook gives you careful explanations, derivations, experiments, exercises and ``a reasonably good guarantee'' that what you learn is true. A press release gives you a fake explanation, conveying only the delusion of understanding, of a result that its writer ``didn't understand'' and that probably will not replicate. \nb{Sumner and colleagues (2014) found exaggerated advice in 40\% of 462 British university press releases, and other kinds of exaggeration in about a third. They did not measure whether the writers understood the results, and the post gives no evidence that they did not.}

Science is built on discoveries built on discoveries, all the way back to Archimedes, who found out why boats float, a fact that makes sense without any other discovery. So a good place to start is the beginning. Do not be embarrassed to read elementary textbooks. ``If you want to pretend to be sophisticated, go find a play to sneer at.'' If you want fun, remember that simplicity is at the core of scientific beauty. Jumping straight to the frontier without the settled science is like trying to climb only the interesting top half of Everest by standing at its base, bending your knees and jumping really hard.

I grant that if 40\% of oncologists thought white socks caused cancer, that would be worth knowing. But science need not be controversial, or recent, to be interesting. Eat only science news, and ``your brain will turn to liquid.''
